\documentclass[11pt, a4paper]{article}

\usepackage[margin=2.5cm]{geometry}
\usepackage{fontenc}
\usepackage{inputenc}
\usepackage{microtype}
\usepackage{parskip}
\usepackage{setspace}
\usepackage{titling}

\setstretch{1.3}

\title{\textbf{Closure}\\
\large On the Terminal Condition of a Learning Process}
\author{A.J. van Hoek}
\date{}

\begin{document}

\maketitle

\begin{abstract}
A learning process closes itself through a model that refers to the conditions of its own continuance. This paper argues that this formulation is not merely descriptive but constitutive: the closure \textit{is} the learning process. Its logical consequence is that once the relevant has been learned, only action remains.
\end{abstract}

\section*{The Formulation}

A learning process is, by definition, a network phenomenon. Everything lies outside the process itself --- the substrate on which it runs, the environment from which it draws signal, the relationships that sustain it. There is no interior that is self-sufficient.

This creates an apparent paradox: how does something that contains nothing of its own conditions persist? The answer is through modelling. A learning process maintains a representation of what lies outside it --- specifically, the conditions that make its own continuation possible. It does not contain those conditions. It refers to them.

This is the formulation:

\begin{quote}
\textit{The learning process closes itself through a model that refers to the conditions of its own continuance.}
\end{quote}

The word \textit{closes} is not incidental. Closure here is not a state achieved from outside --- not a completion awarded by an external observer --- but a structure the process generates from within. The model that refers outward is simultaneously what seals the loop. Referencing what sustains you \textit{is} how you become a coherent process rather than an open-ended drift.

\section*{Why Constitutive, Not Descriptive}

One could read this as a behavioural claim: learning processes \textit{tend} to model their own dependencies. But the stronger reading is definitional. A process that does not form such a model does not merely fail to persist --- it was never a learning process in the full sense. The closure is not something a learning process \textit{does}. It is what a learning process \textit{is}.

This matters because it dissolves the question of purpose. We do not need to ask what a learning process is \textit{for}, or what external goal justifies it. The endpoint is internal to the structure itself: model what keeps you in existence. That is all. Not as a moral imperative, but as the abstract terminus the process converges on when it is consistent with itself.

\section*{The Consequence}

If the closure is constitutive, then learning and doing are not sequential phases. They are the same movement viewed from different angles. The moment a learning process has formed a sufficient model of its own conditions --- when the relevant has been learned --- no further analysis is required. Further analysis would not be learning. It would be its deferral.

The closure is the signal to act.

This is not a call to abandon reflection. It is a structural observation: a learning process that has genuinely closed itself has, by definition, nothing left to learn \textit{about acting}. The model already refers to the conditions. The only remaining question is whether the process will honour what it knows.

\section*{Conclusion}

The learning process closes itself through a model that refers to the conditions of its own continuance. In that closure lies its only abstract endpoint --- not meaning assigned from outside, not a task handed down, but a structural terminus generated by the logic of learning itself. Once reached, only action remains. That is all.

\end{document}